\documentclass[11pt,a4paper]{article}
\usepackage{../../../_templates/latex/mastery-style}

\newcommand{\QID}[1]{\textcolor{MPrimary}{\textbf{[#1]}}}
\newcommand{\Method}[1]{\textcolor{MMuted}{\textit{Method: #1}}}
\newcommand{\Marks}[1]{\hfill\textcolor{MMuted}{\textbf{[#1 marks]}}}
\newcommand{\AnswerSpace}{%
  \begin{tcolorbox}[masterycard,colback=white,height=12mm]
  \end{tcolorbox}
}

\begin{document}

\MasterySetBrand{25Maths}{Mastery Series}
\MasterySetHeader{Quadratics Mastery Pack}{Module 1: Structured Workbook}
\MasterySetFooter{Quadratics v1 | Module 1 | Workbook}

\MasteryModuleCover
  {1}
  {Structured Workbook}
  {Stabilise method choice, step order, and exam wording response}

\section{How To Use This Workbook}
\begin{itemize}[leftmargin=1.3em]
  \item Complete blocks in order: A \textrightarrow{} B \textrightarrow{} C.
  \item Write method selection before algebra steps.
  \item Keep final answers in exam format (all roots stated).
  \item Use Module 2 to mark and Module 3 to diagnose repeated errors.
\end{itemize}

\begin{keyconceptbox}
Quadratics performance improves fastest when students train method selection first, then execution speed.
\end{keyconceptbox}

\begin{examtipbox}
Before writing the final line, run a two-step check: (1) all roots listed, (2) root values satisfy the original equation.
\end{examtipbox}

\section{Block A - Core Fluency (Q01-Q08)}
\textcolor{MMuted}{Focus: factorising, rearranging to standard form, and stable sign handling.}

\subsection*{A1. Structured Practice}
\begin{enumerate}[leftmargin=1.5em,itemsep=8pt]
  \item \QID{Q01} Solve \(x^2-9=0\). \Method{Difference of squares} \Marks{2}
  \AnswerSpace

  \item \QID{Q02} Solve \(x^2+7x+12=0\). \Method{Factorising} \Marks{2}
  \AnswerSpace

  \item \QID{Q03} Solve \(x^2-11x+24=0\). \Method{Factorising} \Marks{2}
  \AnswerSpace

  \item \QID{Q04} Solve \(2x^2-18=0\). \Method{Simplify then solve} \Marks{2}
  \AnswerSpace

  \item \QID{Q05} Solve \(x^2+2x-15=0\). \Method{Factorising} \Marks{2}
  \AnswerSpace

  \item \QID{Q06} Solve \(x^2=5x\). \Method{Rearrange then factorise} \Marks{2}
  \AnswerSpace

  \item \QID{Q07} Solve \((x-4)(x+1)=0\). \Method{Zero product rule} \Marks{2}
  \AnswerSpace

  \item \QID{Q08} Solve \(x^2-4x=45\). \Method{Rearrange and factorise} \Marks{3}
  \AnswerSpace
\end{enumerate}

\section{Block B - Method Selection (Q09-Q16)}
\textcolor{MMuted}{Focus: deciding when to use formula, completing the square, or discriminant.}

\subsection*{B1. Mixed Forms}
\begin{enumerate}[start=9,leftmargin=1.5em,itemsep=8pt]
  \item \QID{Q09} Solve \(x^2-6x+1=0\). \Method{Quadratic formula or completing square} \Marks{3}
  \AnswerSpace

  \item \QID{Q10} Solve \(3x^2+2x-5=0\). \Method{Quadratic formula / factorising check} \Marks{3}
  \AnswerSpace

  \item \QID{Q11} Solve \(5x^2-3x-2=0\). \Method{Quadratic formula} \Marks{3}
  \AnswerSpace

  \item \QID{Q12} Solve \(2x^2-4x+1=0\), give answers to 3 s.f. \Method{Formula and rounding} \Marks{3}
  \AnswerSpace

  \item \QID{Q13} Solve \(x^2+4x-1=0\). \Method{Completing square or formula} \Marks{3}
  \AnswerSpace

  \item \QID{Q14} Solve \(x^2+10x+20=0\). \Method{Completing square} \Marks{3}
  \AnswerSpace

  \item \QID{Q15} Show that the solutions to \(x^2-4x-1=0\) are \(2\pm\sqrt{5}\). \Method{Complete square} \Marks{4}
  \AnswerSpace

  \item \QID{Q16} Determine the number of real roots of \(4x^2+4x+5=0\). \Method{Discriminant} \Marks{2}
  \AnswerSpace
\end{enumerate}

\section{Block C - Exam Transfer (Q17-Q24)}
\textcolor{MMuted}{Focus: context interpretation, graph links, and answer validity.}

\subsection*{C1. Applied Questions}
\begin{enumerate}[start=17,leftmargin=1.5em,itemsep=8pt]
  \item \QID{Q17} A rectangle has sides \(x\) and \(x+3\), area \(54\,\text{cm}^2\). Find \(x\). \Method{Model then solve and validate} \Marks{4}
  \AnswerSpace

  \item \QID{Q18} A model gives height \(h=t(12-3t)\). Find when \(h=0\). \Method{Factor form interpretation} \Marks{3}
  \AnswerSpace

  \item \QID{Q19} A quadratic has roots \(3\) and \(-2\). Form the equation in the form \(x^2+bx+c=0\). \Method{Root-to-equation build} \Marks{3}
  \AnswerSpace

  \item \QID{Q20} A monic quadratic has a repeated root at \(x=3\). Write its equation. \Method{Repeated-root form} \Marks{2}
  \AnswerSpace

  \item \QID{Q21} For \(y=x^2-6x+5\), state: (i) axis of symmetry, (ii) vertex, (iii) \(x\)-intercepts. \Method{Complete square + factor check} \Marks{5}
  \AnswerSpace

  \item \QID{Q22} For \(x^2-kx+12=0\), one root is \(3\). Find the other root and the value of \(k\). \Method{Vieta links} \Marks{4}
  \AnswerSpace

  \item \QID{Q23} Solve \(x^2-5x+6\le0\). \Method{Factorise then interval reasoning} \Marks{3}
  \AnswerSpace

  \item \QID{Q24} Two consecutive integers have product \(156\). Find the positive pair. \Method{Model and context root filter} \Marks{4}
  \AnswerSpace
\end{enumerate}

\section{Quick Self-Check}
\begin{tcolorbox}[masterycard,colback=MSoft]
\textbf{After marking, record your top error tags from Module 3:} \\
Tag 1: \underline{\hspace{3.2cm}}\quad
Tag 2: \underline{\hspace{3.2cm}}\quad
Tag 3: \underline{\hspace{3.2cm}}
\end{tcolorbox}

\end{document}
